% !TEX root = main.tex
\section{Introduction}\label{introduction}

Mental health disorders have surged globally despite economic and technological progress, with loneliness contributing to morbidity risks comparable to obesity {[}1{]}. Population aging, dispersed families, and declining community engagement create unprecedented demand for scalable digital mental health interventions {[}2{]}.

Conversational AI agents for mental health support leveraging large language models (LLMs) have demonstrated capacity for emotionally relevant, context-aware responses {[}8,9{]}, yet a critical limitation persists: many current systems rely on generic or ``one-size-fits-all'' communication strategies and offer only shallow personalization {[}22--24,65,66{]}, while relatively stable personality differences---which shape how individuals process emotions, seek support, and respond to intervention styles {[}4{]}---are seldom incorporated or dynamically adapted {[}65,67{]}. Existing approaches typically use static assessments rather than continuous personality modeling grounded in psychological theory or dynamic adaptation throughout interactions.

To address these limitations, we developed a personality-adaptive conversational framework integrating: (i) real-time Big Five personality trait detection using continuous linguistic analysis, (ii) behavior regulation grounded in the Zurich Model of Social Motivation, explicitly mapping traits to security, arousal, and affiliation domains, and (iii) structured LLM-based evaluation with author-conducted qualitative review. Implemented using GPT-4 and the PROMISE orchestration system {[}7{]}---a model-driven framework that uses state machine modeling to compose prompts from modular parts, enforcing structured prompt/state control and traceable state transitions---the framework performs turn-by-turn personality inference and theory-guided behavioral adaptation throughout multi-turn conversations. The overarching framework was first conceived in prior work {[}68{]} (``a turn-by-turn detection--regulation pipeline for OCEAN-based behavioral adaptation in conversational AI, implemented with GPT-4 and PROMISE orchestration''); the present study extends that conception to a controlled factorial simulation with systematic statistical evaluation.

AI-driven digital interventions in mental health encompass software systems that embed AI techniques to deliver, support, or evaluate mental health services. While conversational AI agents show promise in providing accessible and scalable support, their effectiveness depends on moving beyond generic dialogue capabilities toward personalized interaction strategies grounded in psychological theory. This study contributes to this emerging landscape by demonstrating how relatively stable personality traits can inform adaptive regulation strategies that augment baseline LLM conversational quality with theory-driven personalization, rather than replacing generic competence.

This controlled simulation study tests whether real-time personality-adaptive regulation produces measurable improvements in AI-to-AI simulation compared to non-adaptive baselines across boundary-condition personality profiles. This design provides precise experimental control for quantifying personalization effects, establishes a reproducible methodological testbed, and validates the implementation approach as a foundation for future empirical studies with human participants toward safe, effective, and equitable AI-augmented mental health support. Our results demonstrate near-perfect implementation fidelity (98.3\% detection accuracy, 100\% regulation adherence) and a selective enhancement pattern: dramatic gains in personality-specific need fulfillment under regulation (100\% vs.~8.3\% baseline; Cohen's $d = 4.42$, $p < 0.001$) with no degradation of generic conversational quality, confirming the architectural feasibility of the D--R--E (Detection--Regulation--Evaluation) pipeline.

